\chapter{[MP2] Spezielle Medizinprodukte}
\label{chp:MPG}

\section{Multifunktionsmonitore und Defibrillatoren}

\subsection{Fred}

\subsection{Defigard\texttrademark}

\subsubsection{Defigard 1002}

\subsubsection{Defigard 2002}

\subsubsection{Defigard 6002}

\subsection{Corpuls\texttrademark}



\begin{table}[H]
% \gCaptionof{table}{Corpuls 3: Ausstattungsoptionen.}
\begin{tabulary}{\linewidth}{LL}
\gTabHeading{Merkmal}
&
\gTabHeading{Beschreibung}
\\\addlinespace\midrule\addlinespace
\gTabSub{Hersteller}
&
GS Elektromedizinische Geräte G. Stemple GmbH
\\\addlinespace
\gTabSub{Homepage}
&
\url{http://www.corpuls.com}
\\\addlinespace\bottomrule
\end{tabulary}
\end{table}

\subsubsection{Corpuls 3}

Der Corpuls 3 ist ein Defibrillator/Monitoring-System. Er verlässt den
traditionellen Weg eines klassischen Kompaktgerätes und ist stattdessen modular aufgebaut.

Das System ist teilbar in:

\begin{enumerate}
    \item Monitoreinheit
    \item Patientenbox mit Messmodulen
    \item Defibrillator/Schrittmacher
\end{enumerate}


Die Patientenbox ist mit vielen Messoptionen erhältlich. Die Anschlüsse
befinden sich auf der linken oder rechten Seite der Patientenbox.

\begin{table}[H]

\gCaptionof{table}{Corpuls 3: Ausstattungsoptionen.}
\begin{tabulary}{\linewidth}{LL}
\gTabHeading{Ausstattung}
&
\gTabHeading{Beschreibung}
\\\addlinespace\midrule\addlinespace
\gTabSub{CompactFlash} (Basisausstattung) 
&
Datenspeicherung auf
CompactFlash-Karte, Übergabe der Einsatzdaten u. a. über CompactFlash-Karte möglich
\\\addlinespace
\gTabSub{12-Kanal-EKG} (Basisausstattung)
&
Simultanes 12-Kanal EKG, optional
auch mit Vermessung und Interpretationssoftware
\\\addlinespace
\gTabSub{SpO2}
&
Messung von Sauerstoffsättigung durch Masimo SET-Technologie
\\\addlinespace
\gTabSub{CO2}
&
Kapnometrie, Verfahren auch bei nicht-intubierten Patienten
möglich
\\\addlinespace
\gTabSub{Nicht-invasive Blutdruckmessung} (NIBD)
&
Automatisierte
Messung des Blutdrucks bei Erwachsenen und Neonaten
\\\addlinespace
\gTabSub{2-Kanal Temperatur}  
&
Kerntemperaturmessung invasiv und Hautoberfläche
\\\addlinespace
\gTabSub{Invasive Druckmessung}
&
Bis zu 4 Kanäle IBP, gleichzeitige Messung
des arteriellen, venösen und Gehirndrucks möglich
\\\addlinespace
\gTabSub{Weitere Optionen}
&
Option EKG-Vermessung und EKG-Interpretation, Option Datenübertragung mit
GSM/GPRS \\\addlinespace\bottomrule
\end{tabulary}

\end{table}




\opt{ORG-ASBFLO}{
\ASBFLO{} Es werden folgende Ausstattungsmerkmale eingesetzt:

\begin{compactenum}
    \item CompactFlash
    \item 12-Kanal-EKG
    \item SpO2
    \item CO2
    \item NIBD
\end{compactenum}
}





\subsection{LifePak AED}

\subsection{LifePak\texttrademark Multifunktionsgeräte}

\subsubsection{LifePak 12}

\subsubsection{LifePak 15}

\section{Beatmunggeräte}

\subsection{Weinmann Medumat}

\begin{table}[H]
% \gCaptionof{table}{Corpuls 3: Ausstattungsoptionen.}
\begin{tabulary}{\linewidth}{LL}
\gTabHeading{Merkmal}
&
\gTabHeading{Beschreibung}
\\\addlinespace\midrule\addlinespace
\gTabSub{Hersteller}
&
Weinmann Geräte für Medizin GmbH+Co.KG
\\\addlinespace
\gTabSub{Homepage}
&
\url{http://www.weinmann.de}
\\\addlinespace\bottomrule
\end{tabulary}
\end{table}

\subsubsection{Medumat Compact}

\subsubsection{Medumat Standard und Medumat Standard a}

\begin{figure}[ht]
\includegraphics[width=\linewidth]{./images/geraete/ger-medumat-standard-bedienfeld.jpg}
\gCaptionofDiff{figure}{Bedienfelder Berieselungseinheit \gE{Modul Oxygen} (li.)
und Beatmungsgerät \gE{Medumat{\texttrademark} Standard} (re.).

\gE{Modul Oxygen (li.):} Flowmeter, Ein-Ausschalter und Regelventil. Der
Anschluß links dient der Verbindung der Einheit mit dem Sauerstoffnetz des Fahzeuges.

\gE{Medumat{\texttrademark} Standard (re.):} \gZ{Air-Mix-Schalter,
Ein-Ausschalter, Stellräder für Atemfrequenz, Atemminutenvolumen (!) und Druckbegrenzung,
Beatmungsdruckanzeige, Warnleuchten für Verlegung ("`Stenosis"'), Lösung des
Beatmungsschlauches ("`Disconnection"'), niederen Flaschendruck ("`< 2.7 bar"')
und schwache Batterie.} }
{Bedienfelder Berieselungseinheit \gE{Modul Oxygen}
(li.) und Notfallbeatmungsgerät \gE{Medumat{\texttrademark} Standard} (re.).}
\end{figure}

\subsubsection{Medumat Transport}

\subsection{Draeger Oxylog}

\subsubsection{Oxylog 1000}

\subsubsection{Oxylog 2000}

\subsubsection{Oxylog 3000}